\SourceLeaf{196}\Needspace{10\baselineskip}\begin{center}
\textsc{Remark.}
\end{center}

In a notebook of \textsc{Gauss}, among entries from the
years 1824 to 1828, we find the entry
\begin{center}»\emph{Journal de l'École Normale. Fouriers
Defin. der Ebene. Lacroix I. p. 503.}«\end{center}

Indeed, the Séances des Écoles Normales T. I (Séance du 25
pluviose an III (1795)), Nouvelle Édition, Paris 1800, p. 28, contain remarks
by \textsc{Fourier} in which he proposes new definitions of the straight line and the plane.
\textsc{Fourier} begins, much as \textsc{Bolyai} and
\textsc{Lobatschefskij} do, with the sphere. There seem to be no further publications
by \textsc{Fourier} on the foundations
of geometry.

The Discours préliminaire of the Éléments de Géométrie by \textsc{Lacroix}
(Quatrième édition. Paris, An VIII [1804] p. \textsc{xv}.) contains a reference to
the \emph{Journal des Séances de l'École Normale}. Its omission of
\textsc{Fourier} and mention of \textsc{Laplace} are explained by the fact
that the former had presented his theory in the geometry exercises of the
École Normale conducted by the latter.

\begin{flushright}
\textsc{Stäckel.}
\end{flushright}

\Needspace{10\baselineskip}
\medskip

\begin{center}
\textsc{Gauss} to \textsc{Schumacher}. Göttingen, 11 October 1827.
\end{center}

.\ .\ .\ .\ .\ .\ .\ .\ .\ Who might Mr \textsc{Koch} at No. 83 on the Cremon be? He has
sent me his theory of parallels\HistoricalNote{[*)]}{[*)}{Christian Adolf \textsc{Koch}, \emph{Über Parallellinien.} Ein Versuch, dem Urtheil Sachkundiger gewidmet. Hamburg 1827. $8^{0}$. 12 pp.]} which is of course without merit. But
it is a rare thing that he at once acknowledged the flaw I pointed out,
as he did my comments on the equally unsuccessful attempts that he later sent me
in writing. .\ .\ .\ .\ .

